\clearpage\pagestyle{otMain}\setlength{\otBlocEspacement}{2.5mm}\fontsize{11}{15}\selectfont\corrigesChapitre{8}{Fractions décimales et virgule}
\begin{blocLivre}{cor-8-1}\begin{corrige}{1}{Des dixièmes à lire}\demarche{Chercher les unités entières dans les dixièmes.}
« $\res{\text{sept dixièmes}}$ » et « $\res{\text{vingt-trois dixièmes}}$ ». Comme $\frac{20}{10}=2$ : $\res{\frac{23}{10}=2+\frac3{10}}$.
\end{corrige}\end{blocLivre}
\begin{blocLivre}{cor-8-2}\begin{corrige}{2}{Sur la demi-droite}\demarche{Compter les pas de un dixième.}
A est 4 pas après 2 : $\res{A\left(\frac{24}{10}\right)}$. B est 8 pas après 2 : $\res{B\left(\frac{28}{10}\right)}$.
\label{fig-cor-8-2}\figureGeo[corrige]{dec-2}
\end{corrige}\end{blocLivre}
\begin{blocLivre}{cor-8-3}\begin{corrige}{3}{Unités, dixièmes, centièmes}\demarche{Un dixième vaut dix centièmes.}
Une unité vaut $\res{100}$ centièmes. Trois dixièmes font $\res{30}$ centièmes. Cinquante centièmes font $\res5$ dixièmes.
\end{corrige}\end{blocLivre}
\begin{blocLivre}{cor-8-4}\begin{corrige}{4}{Vers la virgule}\demarche{Placer chaque chiffre à son rang.}
La partie entière est 35, puis 7 dixièmes et 8 centièmes : $\res{35{,}78}$. Le second nombre n’a pas d’unité entière : $\res{0{,}42}$.
\end{corrige}\end{blocLivre}
\begin{blocLivre}{cor-8-5}\begin{corrige}{5}{Depuis la virgule}\demarche{Lire chaque chiffre avec son rang.}
$6{,}09=\res{6+\frac9{100}}$ : le zéro des dixièmes ne s’écrit pas dans la somme. $14{,}3=\res{14+\frac3{10}}$.
\end{corrige}\end{blocLivre}
\begin{blocLivre}{cor-8-6}\begin{corrige}{6}{Fractions en virgule}\demarche{Extraire d’abord la partie entière.}
$\frac{45}{10}=4+\frac5{10}=\res{4{,}5}$. $\frac{307}{100}=3+\frac7{100}=\res{3{,}07}$. $\frac8{100}=\res{0{,}08}$.
\end{corrige}\end{blocLivre}
\begin{blocLivre}{cor-8-7}\begin{corrige}{7}{Le rang de chaque chiffre}\demarche{Repérer la virgule puis compter les rangs.}
Le chiffre des dixièmes est $\res3$, celui des centièmes $\res6$. Le 3 est au rang des dixièmes : il vaut $\res{\frac3{10}}$.
\end{corrige}\end{blocLivre}
\begin{blocLivre}{cor-8-8}\begin{corrige}{8}{Deux écritures, un nombre}\demarche{Passer par la décomposition en rangs.}
Quatre unités et sept centièmes : $\res{\frac{407}{100}=4{,}07}$. Douze dixièmes : $\res{\frac{12}{10}=1{,}2}$ , soit une unité et deux dixièmes.
\end{corrige}\end{blocLivre}
\begin{blocLivre}{cor-8-9}\begin{corrige}{9}{Avec la monnaie}\demarche{Compter les centièmes d’euro.}
$3+\frac{45}{100}=\res{3{,}45}$ euro. $275$ centimes font $\frac{275}{100}=\res{2{,}75}$ euros. Huit centimes : $\res{0{,}08}$ euro.
\end{corrige}\end{blocLivre}
\begin{blocLivre}{cor-8-10}\begin{corrige}{10}{Des zéros qui comptent}\demarche{Associer chaque chiffre à son rang.}
Cinq centièmes : $\res{0{,}05}$. Cinq dixièmes : $\res{0{,}5}$. Ils sont $\res{\text{différents}}$ : $\frac5{10}=\frac{50}{100}$, dix fois plus que $\frac5{100}$.
\end{corrige}\end{blocLivre}
\begin{blocLivre}{cor-8-11}\begin{corrige}{11}{La grille aux cent cases}\demarche{Lire les centièmes directement sur la grille.}
$\res{\frac{37}{100}=0{,}37}$ de la grille. Avec 3 cases de plus : $\frac{40}{100}=\res{0{,}40}$, c’est-à-dire $0{,}4$ : quatre dixièmes, soit quatre colonnes entières.
\label{fig-cor-8-11}\figureGeo[corrige]{dec-13}
\end{corrige}\end{blocLivre}
\begin{blocLivre}{cor-8-12}\begin{corrige}{12}{Le compteur en dixièmes}\demarche{Compter de dixième en dixième avec les échanges.}
$45{,}8$ puis $\res{45{,}9}$ puis $\res{46}$ : neuf dixièmes plus un dixième font dix dixièmes, une unité entière. Le compteur ne passe pas par « $45{,}10$ ».
\end{corrige}\end{blocLivre}
\begin{blocLivre}{cor-8-13}\begin{corrige}{13}{Les sauts mesurés}\demarche{Décomposer rang par rang, sans oublier les zéros.}
$3{,}45=3+\frac4{10}+\frac5{100}$ ; $3{,}5=3+\frac5{10}$ ; $3{,}09=3+\frac9{100}$. Les chiffres des dixièmes sont $\res4$, $\res5$ et $\res0$.
\end{corrige}\end{blocLivre}
\begin{blocLivre}{cor-8-14}\begin{corrige}{14}{Fractions bien connues}\demarche{Convertir en centièmes avant de coder.}
$\frac12=\frac{50}{100}=\res{0{,}5}$ ; $\frac14=\frac{25}{100}=\res{0{,}25}$ ; $\frac34=\frac{75}{100}=\res{0{,}75}$. Ces écritures servent très souvent en calcul mental.
\end{corrige}\end{blocLivre}
\begin{blocLivre}{cor-8-15}\begin{corrige}{15}{Vrai ou faux ?}\demarche{Traduire chaque écriture en fractions décimales.}
Milo se trompe : $0{,}7=\frac7{10}=\frac{70}{100}$, soit $\res{70}$ centièmes. Sara a raison : $2{,}30=2+\frac{30}{100}=2+\frac3{10}=\res{2{,}3}$. Un zéro final à droite ne change pas le nombre ; un zéro entre deux chiffres, si.
\end{corrige}\end{blocLivre}
